Fix \(P\in\mathfrak P_{\mathcal D}^{\mathrm{obs}}\), \(d\in\mathcal D\), and \(0<\alpha<1\), and put
\[
 x_E=\frac{\alpha B}{c_{E,d}},
 \qquad
 x_O=\frac{(1-\alpha)B}{c_{O,d}}.
\]
The stated lower bound on \(B\) implies \(x_E,x_O\ge2\).  For every \(x\ge2\),
\[
 0\le \frac1{\lfloor x\rfloor}-\frac1x
 =\frac{x-\lfloor x\rfloor}{x\lfloor x\rfloor}
 \le\frac2{x^2},                                           \tag{1}
\]
because \(0\le x-\lfloor x\rfloor<1\) and \(\lfloor x\rfloor\ge x/2\).  Apply (1) first to \(x_E\) and then to \(x_O\).  The definitions of the two floor counts, the design-value profile, and the residual then give the exact identity
\[
 B\left\{\frac{V_{E,d}}{n_{E,B}(d,\alpha)}
 +\frac{V_{O,d}}{n_{O,B}(d,\alpha)}\right\}
 =\mathcal V_d(\alpha;P)+\mathfrak r_{d,B}(\alpha;P)
\]
and the bound
\[
 0\le\mathfrak r_{d,B}(\alpha;P)
 \le\frac2B\left\{
 \frac{V_{E,d}c_{E,d}^2}{\alpha^2}
 +\frac{V_{O,d}c_{O,d}^2}{(1-\alpha)^2}\right\}.       \tag{2}
\]
For fixed \(d\) and \(\alpha\), the right side of (2) tends to zero.  Lemma~\(\mathrm{lem{:}designwise\text{-}gradient}\) identifies the expression inside braces before multiplication by \(B\) as the floor-count variance of the two-source gradient of \(\theta_d(P)\).  This proves the claimed convergence.

Set
\[
 a_d=c_{E,d}V_{E,d},
 \qquad
 b_d=c_{O,d}V_{O,d}.
\]
Positive costs and nondegeneracy imply \(a_d,b_d>0\).  Direct algebra gives, for \(0<\alpha<1\),
\[
 \frac{a_d}{\alpha}+\frac{b_d}{1-\alpha}
 -(\sqrt{a_d}+\sqrt{b_d})^2
 =\frac{\{(1-\alpha)\sqrt{a_d}-\alpha\sqrt{b_d}\}^2}
 {\alpha(1-\alpha)}\ge0.                                  \tag{3}
\]
Equality holds if and only if
\[
 \alpha
 =\frac{\sqrt{a_d}}{\sqrt{a_d}+\sqrt{b_d}}
 =\frac{\sqrt{c_{E,d}V_{E,d}}}
 {\sqrt{c_{E,d}V_{E,d}}+\sqrt{c_{O,d}V_{O,d}}}.
\]
Thus the minimizer is unique and equals \(\alpha_d^*(P)\), and
\[
 \mathcal V_d^*(P)
 =(\sqrt{c_{E,d}V_{E,d}}+\sqrt{c_{O,d}V_{O,d}})^2.
\]
If \(P\in\mathfrak P_{\mathcal D}^{\mathrm{id}}\), the identity \(\theta_d(P)=\tau\) holds for every admissible design, so the displayed profile is the gradient-variance profile of the common causal target along that subclass.  If, in addition, this designwise gradient is the canonical gradient for \(\tau\), its variance is by definition the local semiparametric efficiency lower bound.  The fixed-system Imbens--Kallus--Mao--Wang premises recorded separately are one setting in which that additional causal canonical-gradient conclusion may be checked; none of it is needed for (1)--(3).